% =====================================================================
%  ablation.tex  --  PHASE VI, Task 13E
%
%  The component ablation: what each stage of the pipeline contributes,
%  measured in the metric that stage is responsible for.
%
%  Intended placement: Section 8 (\label{sec:ablation}), after the
%  explainability section.
%
%  Drop-in usage: % =====================================================================
%  ablation.tex  --  PHASE VI, Task 13E
%
%  The component ablation: what each stage of the pipeline contributes,
%  measured in the metric that stage is responsible for.
%
%  Intended placement: Section 8 (\label{sec:ablation}), after the
%  explainability section.
%
%  Drop-in usage: % =====================================================================
%  ablation.tex  --  PHASE VI, Task 13E
%
%  The component ablation: what each stage of the pipeline contributes,
%  measured in the metric that stage is responsible for.
%
%  Intended placement: Section 8 (\label{sec:ablation}), after the
%  explainability section.
%
%  Drop-in usage: % =====================================================================
%  ablation.tex  --  PHASE VI, Task 13E
%
%  The component ablation: what each stage of the pipeline contributes,
%  measured in the metric that stage is responsible for.
%
%  Intended placement: Section 8 (\label{sec:ablation}), after the
%  explainability section.
%
%  Drop-in usage: \input{ablation}
%                 \input{../outputs/phase6/task13e/pelletizer-I/ablation_table}
%  Figures are referenced by bare name; add to the preamble
%      \graphicspath{{../outputs/phase6/figures/}}
%  The table needs \usepackage{booktabs}.
%  Numbers below are generated by  python run_phase6.py --only task13e
%  and stored in outputs/phase6/task13e/<machine>/.
%  The table itself is GENERATED -- do not edit ablation_table.tex by
%  hand; edit the script that writes it.
%
%  FORWARD REFERENCES: sec:forecasting, sec:significance (see the note
%  at the head of explainability.tex).
% =====================================================================

\section{Component Ablation}
\label{sec:ablation}

A framework assembled from five components invites one question above all
others: does each of them earn its place? Table~\ref{tab:ablation} answers it
by removing one component at a time and reporting what the removal costs.
Every cell is read from the stored output of the experiment that produced the
corresponding headline result, so the ablation and the claim it qualifies come
from the same fitted models and the same held-out data; nothing here is
refitted for the purpose.

\subsection{Why the table is blocked rather than cumulative}
\label{sec:ablation:structure}

The three stages are measured on three different units --- readings, windows
and factory days --- so a single column running the height of the table would
be meaningless. Two structural points follow, and both matter for reading it.

First, the state-layer metrics are held constant below the first block rather
than repeated: every forecast and decision configuration runs on the
minimum-dwell labelling, and printing its flicker rate nine times would assert
nine times what is asserted once.

Second, and less obviously, \textbf{the two forecasters in this framework are
different models and the table does not pretend otherwise}. The window
forecaster of Section~\ref{sec:forecasting} predicts the state at each step of
a 300\,s horizon and is scored on windows; the decision layer does not consume
it. What the decision layer consumes is the episode-duration forecaster of
Section~\ref{sec:decision:policies}, whose error is measured in seconds of
remaining idle time ($8\,954\pm156$\,s over ten seeds). Carrying the window
model's $F_1$ into the decision rows would attribute the economic result to a
model that plays no part in producing it.

\subsection{The state layer}
\label{sec:ablation:state}

Flicker --- the fraction of state runs shorter than the 10\,s physical
plausibility floor --- falls from $66.9\%$ under independent GMM assignment to
$44.5\%$ once Viterbi decoding imposes a temporal prior, and to $0.00\%$ once
an explicit minimum-dwell constraint is added. The median dwell rises from
3\,s to 13\,s to 92.5\,s across the same three configurations.

Two readings of this are worth separating. \textbf{The HMM contributes, but it
does not solve the problem}: a transition prior can only make short runs
unlikely, and at $10^{6}$ readings per record the likelihood term overwhelms
any prior that leaves the model able to change state at all. It takes an
explicit dwell floor to remove flicker outright, which is what
Section~\ref{sec:method:ghmm} argues on theoretical grounds and what this row
of the table measures.

\textbf{Neither component changes the physics compliance}, and the table says
so rather than hiding it: all three configurations pass 5/5, and the schedule
score moves by $0.06$ percentage points across the block. This is the expected
result and not a null one. The dwell constraint changes \emph{when} the machine
is said to change state, not \emph{which} cluster is standby; a table showing
it improving the physics score would indicate a leak between the two, not a
better model. For scale, the three clustering baselines outside the ablation
reach flicker rates of $62.5\%$ (DBSCAN), $78.2\%$ (K-Means) and $98.2\%$
(spectral clustering), the last of which also fails a physics check.

\subsection{The forecast layer}
\label{sec:ablation:forecast}

The block is anchored on an untrained reference --- the horizon continues the
last observed state --- because a comparison among trained models cannot show
that training helped.

The sequence model does not survive that comparison. The Seq2Seq LSTM reaches
STANDBY $F_1 = 0.561\pm0.064$ against persistence's $0.681$, and $20.90$\,s of
per-window MAE against $12.95$\,s: worse than not forecasting, by
$7.95$\,s, at $p_{\text{Holm}} = 8.6\times10^{-43}$
(Section~\ref{sec:significance}). The gradient-boosted tree is the one trained
forecaster that clears the reference, at $0.689\pm0.003$ and $11.50$\,s ---
$1.45$\,s better than persistence, $p_{\text{Holm}} = 7\times10^{-4}$, and
above the pre-registered 1\,s practical threshold.

The standard deviations in the block carry a second finding. The tree's
spread over 15 seeds is $0.003$ of $F_1$; the sequence model's over 5 seeds is
$0.064$, a factor of twenty-five. The swing between a good and a bad seed of
one neural architecture is larger than the entire spread between
architectures, which is why the forecasting comparison is reported as
mean\,$\pm$\,sd over seeds and not as a single-seed ranking.

\subsection{The decision layer}
\label{sec:ablation:decision}

Every decision row carries a 95\% day-block bootstrap interval over 18 factory
days, and the intervals are the point of the block. Two readings of them must
be kept apart. Each row \emph{on its own} is an interval against the status
quo, and every one of those spans zero except the oracle's: the proposed
policy's $+2.9$ USD/yr sits inside $[-55.9,+46.2]$, so no single row
demonstrates a saving. A \emph{paired contrast} between two rows is a different
and much sharper quantity, because every policy is scored on the same bootstrap
resample and the pairing removes the day-to-day variance that makes the
individual intervals wide. The contrasts quoted below are of the second kind,
and one of them survives.

Three components are separable, and their contributions differ by an order of
magnitude.

\textbf{Targeting the right functional of the predictive distribution is worth
more than anything else in the layer.} Replacing the conditional-mean
forecaster with a log-scale one --- the natural-looking choice, which estimates
the conditional median --- costs $31.5$ USD/yr, from $+2.9$ to $-28.6$. The
idle-duration distribution is heavy enough that its conditional mean exceeds
its conditional median by more than an order of magnitude at episode onset, and
the saving from de-energising is affine in the remaining duration, so the mean
is the sufficient statistic and a policy driven by the median waits far too
long to act. This is the single largest ablation effect in the table.

\textbf{The chance constraint is worth $17.8$ USD/yr}, from $-15.0$ without it
to $+2.9$ with it, and it is what moves the policy's point estimate across
zero. It also enforces the constraint it was introduced for. On the Phase~I--III
decision run, the only one for which both variants' violation counts are
recorded, the unconstrained policy leaves 19 minimum-off violations against the
proposed policy's 2; under the minimum-dwell labelling the proposed policy's
count is again 2, with no variation over ten seeds.

\textbf{The optimisation layer beats the fixed threshold it was built to
replace, and this is the one decision-layer contrast that is significant.} The
proposed policy returns $+68.9$ USD/yr more than the static break-even rule
($p = 0.001$, interval excluding zero, and above the 5 USD/yr economic
threshold). It does \emph{not} beat forecast-free ski-rental
($+2.8$, $p = 0.729$) and cannot be shown to beat the plant's own status quo
($+2.9$, $p = 0.828$).

\textbf{The conclusion the table supports is therefore narrow and should be
stated narrowly}: adaptive thresholds beat a fixed one; forecasting does not
demonstrably beat not forecasting on this machine and this sample. The oracle
row bounds what was available --- $+37.6$ USD/yr with an interval excluding
zero --- so the head-room is real and largely uncaptured, and the binding
constraint on any stronger claim is 18 factory days, not the estimator.

\begin{figure}[t]
  \centering
  \includegraphics[width=\linewidth]{fig_task13e_ablation}
  \caption{The ablation of Table~\ref{tab:ablation} in three panels, one per
  stage, because the three stages are measured in three units. (A)~Flicker
  under each temporal component, pooled over all states and binarised to
  STANDBY-versus-rest, which is the view the decision layer is sensitive to.
  (B)~Per-window duration error with the seed spread; the dashed line is the
  untrained reference. (C)~Annual saving with 95\% day-block bootstrap
  intervals; hollow markers mark intervals that include zero.}
  \label{fig:ablation}
\end{figure}

%                 \input{../outputs/phase6/task13e/pelletizer-I/ablation_table}
%  Figures are referenced by bare name; add to the preamble
%      \graphicspath{{../outputs/phase6/figures/}}
%  The table needs \usepackage{booktabs}.
%  Numbers below are generated by  python run_phase6.py --only task13e
%  and stored in outputs/phase6/task13e/<machine>/.
%  The table itself is GENERATED -- do not edit ablation_table.tex by
%  hand; edit the script that writes it.
%
%  FORWARD REFERENCES: sec:forecasting, sec:significance (see the note
%  at the head of explainability.tex).
% =====================================================================

\section{Component Ablation}
\label{sec:ablation}

A framework assembled from five components invites one question above all
others: does each of them earn its place? Table~\ref{tab:ablation} answers it
by removing one component at a time and reporting what the removal costs.
Every cell is read from the stored output of the experiment that produced the
corresponding headline result, so the ablation and the claim it qualifies come
from the same fitted models and the same held-out data; nothing here is
refitted for the purpose.

\subsection{Why the table is blocked rather than cumulative}
\label{sec:ablation:structure}

The three stages are measured on three different units --- readings, windows
and factory days --- so a single column running the height of the table would
be meaningless. Two structural points follow, and both matter for reading it.

First, the state-layer metrics are held constant below the first block rather
than repeated: every forecast and decision configuration runs on the
minimum-dwell labelling, and printing its flicker rate nine times would assert
nine times what is asserted once.

Second, and less obviously, \textbf{the two forecasters in this framework are
different models and the table does not pretend otherwise}. The window
forecaster of Section~\ref{sec:forecasting} predicts the state at each step of
a 300\,s horizon and is scored on windows; the decision layer does not consume
it. What the decision layer consumes is the episode-duration forecaster of
Section~\ref{sec:decision:policies}, whose error is measured in seconds of
remaining idle time ($8\,954\pm156$\,s over ten seeds). Carrying the window
model's $F_1$ into the decision rows would attribute the economic result to a
model that plays no part in producing it.

\subsection{The state layer}
\label{sec:ablation:state}

Flicker --- the fraction of state runs shorter than the 10\,s physical
plausibility floor --- falls from $66.9\%$ under independent GMM assignment to
$44.5\%$ once Viterbi decoding imposes a temporal prior, and to $0.00\%$ once
an explicit minimum-dwell constraint is added. The median dwell rises from
3\,s to 13\,s to 92.5\,s across the same three configurations.

Two readings of this are worth separating. \textbf{The HMM contributes, but it
does not solve the problem}: a transition prior can only make short runs
unlikely, and at $10^{6}$ readings per record the likelihood term overwhelms
any prior that leaves the model able to change state at all. It takes an
explicit dwell floor to remove flicker outright, which is what
Section~\ref{sec:method:ghmm} argues on theoretical grounds and what this row
of the table measures.

\textbf{Neither component changes the physics compliance}, and the table says
so rather than hiding it: all three configurations pass 5/5, and the schedule
score moves by $0.06$ percentage points across the block. This is the expected
result and not a null one. The dwell constraint changes \emph{when} the machine
is said to change state, not \emph{which} cluster is standby; a table showing
it improving the physics score would indicate a leak between the two, not a
better model. For scale, the three clustering baselines outside the ablation
reach flicker rates of $62.5\%$ (DBSCAN), $78.2\%$ (K-Means) and $98.2\%$
(spectral clustering), the last of which also fails a physics check.

\subsection{The forecast layer}
\label{sec:ablation:forecast}

The block is anchored on an untrained reference --- the horizon continues the
last observed state --- because a comparison among trained models cannot show
that training helped.

The sequence model does not survive that comparison. The Seq2Seq LSTM reaches
STANDBY $F_1 = 0.561\pm0.064$ against persistence's $0.681$, and $20.90$\,s of
per-window MAE against $12.95$\,s: worse than not forecasting, by
$7.95$\,s, at $p_{\text{Holm}} = 8.6\times10^{-43}$
(Section~\ref{sec:significance}). The gradient-boosted tree is the one trained
forecaster that clears the reference, at $0.689\pm0.003$ and $11.50$\,s ---
$1.45$\,s better than persistence, $p_{\text{Holm}} = 7\times10^{-4}$, and
above the pre-registered 1\,s practical threshold.

The standard deviations in the block carry a second finding. The tree's
spread over 15 seeds is $0.003$ of $F_1$; the sequence model's over 5 seeds is
$0.064$, a factor of twenty-five. The swing between a good and a bad seed of
one neural architecture is larger than the entire spread between
architectures, which is why the forecasting comparison is reported as
mean\,$\pm$\,sd over seeds and not as a single-seed ranking.

\subsection{The decision layer}
\label{sec:ablation:decision}

Every decision row carries a 95\% day-block bootstrap interval over 18 factory
days, and the intervals are the point of the block. Two readings of them must
be kept apart. Each row \emph{on its own} is an interval against the status
quo, and every one of those spans zero except the oracle's: the proposed
policy's $+2.9$ USD/yr sits inside $[-55.9,+46.2]$, so no single row
demonstrates a saving. A \emph{paired contrast} between two rows is a different
and much sharper quantity, because every policy is scored on the same bootstrap
resample and the pairing removes the day-to-day variance that makes the
individual intervals wide. The contrasts quoted below are of the second kind,
and one of them survives.

Three components are separable, and their contributions differ by an order of
magnitude.

\textbf{Targeting the right functional of the predictive distribution is worth
more than anything else in the layer.} Replacing the conditional-mean
forecaster with a log-scale one --- the natural-looking choice, which estimates
the conditional median --- costs $31.5$ USD/yr, from $+2.9$ to $-28.6$. The
idle-duration distribution is heavy enough that its conditional mean exceeds
its conditional median by more than an order of magnitude at episode onset, and
the saving from de-energising is affine in the remaining duration, so the mean
is the sufficient statistic and a policy driven by the median waits far too
long to act. This is the single largest ablation effect in the table.

\textbf{The chance constraint is worth $17.8$ USD/yr}, from $-15.0$ without it
to $+2.9$ with it, and it is what moves the policy's point estimate across
zero. It also enforces the constraint it was introduced for. On the Phase~I--III
decision run, the only one for which both variants' violation counts are
recorded, the unconstrained policy leaves 19 minimum-off violations against the
proposed policy's 2; under the minimum-dwell labelling the proposed policy's
count is again 2, with no variation over ten seeds.

\textbf{The optimisation layer beats the fixed threshold it was built to
replace, and this is the one decision-layer contrast that is significant.} The
proposed policy returns $+68.9$ USD/yr more than the static break-even rule
($p = 0.001$, interval excluding zero, and above the 5 USD/yr economic
threshold). It does \emph{not} beat forecast-free ski-rental
($+2.8$, $p = 0.729$) and cannot be shown to beat the plant's own status quo
($+2.9$, $p = 0.828$).

\textbf{The conclusion the table supports is therefore narrow and should be
stated narrowly}: adaptive thresholds beat a fixed one; forecasting does not
demonstrably beat not forecasting on this machine and this sample. The oracle
row bounds what was available --- $+37.6$ USD/yr with an interval excluding
zero --- so the head-room is real and largely uncaptured, and the binding
constraint on any stronger claim is 18 factory days, not the estimator.

\begin{figure}[t]
  \centering
  \includegraphics[width=\linewidth]{fig_task13e_ablation}
  \caption{The ablation of Table~\ref{tab:ablation} in three panels, one per
  stage, because the three stages are measured in three units. (A)~Flicker
  under each temporal component, pooled over all states and binarised to
  STANDBY-versus-rest, which is the view the decision layer is sensitive to.
  (B)~Per-window duration error with the seed spread; the dashed line is the
  untrained reference. (C)~Annual saving with 95\% day-block bootstrap
  intervals; hollow markers mark intervals that include zero.}
  \label{fig:ablation}
\end{figure}

%                 \input{../outputs/phase6/task13e/pelletizer-I/ablation_table}
%  Figures are referenced by bare name; add to the preamble
%      \graphicspath{{../outputs/phase6/figures/}}
%  The table needs \usepackage{booktabs}.
%  Numbers below are generated by  python run_phase6.py --only task13e
%  and stored in outputs/phase6/task13e/<machine>/.
%  The table itself is GENERATED -- do not edit ablation_table.tex by
%  hand; edit the script that writes it.
%
%  FORWARD REFERENCES: sec:forecasting, sec:significance (see the note
%  at the head of explainability.tex).
% =====================================================================

\section{Component Ablation}
\label{sec:ablation}

A framework assembled from five components invites one question above all
others: does each of them earn its place? Table~\ref{tab:ablation} answers it
by removing one component at a time and reporting what the removal costs.
Every cell is read from the stored output of the experiment that produced the
corresponding headline result, so the ablation and the claim it qualifies come
from the same fitted models and the same held-out data; nothing here is
refitted for the purpose.

\subsection{Why the table is blocked rather than cumulative}
\label{sec:ablation:structure}

The three stages are measured on three different units --- readings, windows
and factory days --- so a single column running the height of the table would
be meaningless. Two structural points follow, and both matter for reading it.

First, the state-layer metrics are held constant below the first block rather
than repeated: every forecast and decision configuration runs on the
minimum-dwell labelling, and printing its flicker rate nine times would assert
nine times what is asserted once.

Second, and less obviously, \textbf{the two forecasters in this framework are
different models and the table does not pretend otherwise}. The window
forecaster of Section~\ref{sec:forecasting} predicts the state at each step of
a 300\,s horizon and is scored on windows; the decision layer does not consume
it. What the decision layer consumes is the episode-duration forecaster of
Section~\ref{sec:decision:policies}, whose error is measured in seconds of
remaining idle time ($8\,954\pm156$\,s over ten seeds). Carrying the window
model's $F_1$ into the decision rows would attribute the economic result to a
model that plays no part in producing it.

\subsection{The state layer}
\label{sec:ablation:state}

Flicker --- the fraction of state runs shorter than the 10\,s physical
plausibility floor --- falls from $66.9\%$ under independent GMM assignment to
$44.5\%$ once Viterbi decoding imposes a temporal prior, and to $0.00\%$ once
an explicit minimum-dwell constraint is added. The median dwell rises from
3\,s to 13\,s to 92.5\,s across the same three configurations.

Two readings of this are worth separating. \textbf{The HMM contributes, but it
does not solve the problem}: a transition prior can only make short runs
unlikely, and at $10^{6}$ readings per record the likelihood term overwhelms
any prior that leaves the model able to change state at all. It takes an
explicit dwell floor to remove flicker outright, which is what
Section~\ref{sec:method:ghmm} argues on theoretical grounds and what this row
of the table measures.

\textbf{Neither component changes the physics compliance}, and the table says
so rather than hiding it: all three configurations pass 5/5, and the schedule
score moves by $0.06$ percentage points across the block. This is the expected
result and not a null one. The dwell constraint changes \emph{when} the machine
is said to change state, not \emph{which} cluster is standby; a table showing
it improving the physics score would indicate a leak between the two, not a
better model. For scale, the three clustering baselines outside the ablation
reach flicker rates of $62.5\%$ (DBSCAN), $78.2\%$ (K-Means) and $98.2\%$
(spectral clustering), the last of which also fails a physics check.

\subsection{The forecast layer}
\label{sec:ablation:forecast}

The block is anchored on an untrained reference --- the horizon continues the
last observed state --- because a comparison among trained models cannot show
that training helped.

The sequence model does not survive that comparison. The Seq2Seq LSTM reaches
STANDBY $F_1 = 0.561\pm0.064$ against persistence's $0.681$, and $20.90$\,s of
per-window MAE against $12.95$\,s: worse than not forecasting, by
$7.95$\,s, at $p_{\text{Holm}} = 8.6\times10^{-43}$
(Section~\ref{sec:significance}). The gradient-boosted tree is the one trained
forecaster that clears the reference, at $0.689\pm0.003$ and $11.50$\,s ---
$1.45$\,s better than persistence, $p_{\text{Holm}} = 7\times10^{-4}$, and
above the pre-registered 1\,s practical threshold.

The standard deviations in the block carry a second finding. The tree's
spread over 15 seeds is $0.003$ of $F_1$; the sequence model's over 5 seeds is
$0.064$, a factor of twenty-five. The swing between a good and a bad seed of
one neural architecture is larger than the entire spread between
architectures, which is why the forecasting comparison is reported as
mean\,$\pm$\,sd over seeds and not as a single-seed ranking.

\subsection{The decision layer}
\label{sec:ablation:decision}

Every decision row carries a 95\% day-block bootstrap interval over 18 factory
days, and the intervals are the point of the block. Two readings of them must
be kept apart. Each row \emph{on its own} is an interval against the status
quo, and every one of those spans zero except the oracle's: the proposed
policy's $+2.9$ USD/yr sits inside $[-55.9,+46.2]$, so no single row
demonstrates a saving. A \emph{paired contrast} between two rows is a different
and much sharper quantity, because every policy is scored on the same bootstrap
resample and the pairing removes the day-to-day variance that makes the
individual intervals wide. The contrasts quoted below are of the second kind,
and one of them survives.

Three components are separable, and their contributions differ by an order of
magnitude.

\textbf{Targeting the right functional of the predictive distribution is worth
more than anything else in the layer.} Replacing the conditional-mean
forecaster with a log-scale one --- the natural-looking choice, which estimates
the conditional median --- costs $31.5$ USD/yr, from $+2.9$ to $-28.6$. The
idle-duration distribution is heavy enough that its conditional mean exceeds
its conditional median by more than an order of magnitude at episode onset, and
the saving from de-energising is affine in the remaining duration, so the mean
is the sufficient statistic and a policy driven by the median waits far too
long to act. This is the single largest ablation effect in the table.

\textbf{The chance constraint is worth $17.8$ USD/yr}, from $-15.0$ without it
to $+2.9$ with it, and it is what moves the policy's point estimate across
zero. It also enforces the constraint it was introduced for. On the Phase~I--III
decision run, the only one for which both variants' violation counts are
recorded, the unconstrained policy leaves 19 minimum-off violations against the
proposed policy's 2; under the minimum-dwell labelling the proposed policy's
count is again 2, with no variation over ten seeds.

\textbf{The optimisation layer beats the fixed threshold it was built to
replace, and this is the one decision-layer contrast that is significant.} The
proposed policy returns $+68.9$ USD/yr more than the static break-even rule
($p = 0.001$, interval excluding zero, and above the 5 USD/yr economic
threshold). It does \emph{not} beat forecast-free ski-rental
($+2.8$, $p = 0.729$) and cannot be shown to beat the plant's own status quo
($+2.9$, $p = 0.828$).

\textbf{The conclusion the table supports is therefore narrow and should be
stated narrowly}: adaptive thresholds beat a fixed one; forecasting does not
demonstrably beat not forecasting on this machine and this sample. The oracle
row bounds what was available --- $+37.6$ USD/yr with an interval excluding
zero --- so the head-room is real and largely uncaptured, and the binding
constraint on any stronger claim is 18 factory days, not the estimator.

\begin{figure}[t]
  \centering
  \includegraphics[width=\linewidth]{fig_task13e_ablation}
  \caption{The ablation of Table~\ref{tab:ablation} in three panels, one per
  stage, because the three stages are measured in three units. (A)~Flicker
  under each temporal component, pooled over all states and binarised to
  STANDBY-versus-rest, which is the view the decision layer is sensitive to.
  (B)~Per-window duration error with the seed spread; the dashed line is the
  untrained reference. (C)~Annual saving with 95\% day-block bootstrap
  intervals; hollow markers mark intervals that include zero.}
  \label{fig:ablation}
\end{figure}

%                 \input{../outputs/phase6/task13e/pelletizer-I/ablation_table}
%  Figures are referenced by bare name; add to the preamble
%      \graphicspath{{../outputs/phase6/figures/}}
%  The table needs \usepackage{booktabs}.
%  Numbers below are generated by  python run_phase6.py --only task13e
%  and stored in outputs/phase6/task13e/<machine>/.
%  The table itself is GENERATED -- do not edit ablation_table.tex by
%  hand; edit the script that writes it.
%
%  FORWARD REFERENCES: sec:forecasting, sec:significance (see the note
%  at the head of explainability.tex).
% =====================================================================

\section{Component Ablation}
\label{sec:ablation}

A framework assembled from five components invites one question above all
others: does each of them earn its place? Table~\ref{tab:ablation} answers it
by removing one component at a time and reporting what the removal costs.
Every cell is read from the stored output of the experiment that produced the
corresponding headline result, so the ablation and the claim it qualifies come
from the same fitted models and the same held-out data; nothing here is
refitted for the purpose.

\subsection{Why the table is blocked rather than cumulative}
\label{sec:ablation:structure}

The three stages are measured on three different units --- readings, windows
and factory days --- so a single column running the height of the table would
be meaningless. Two structural points follow, and both matter for reading it.

First, the state-layer metrics are held constant below the first block rather
than repeated: every forecast and decision configuration runs on the
minimum-dwell labelling, and printing its flicker rate nine times would assert
nine times what is asserted once.

Second, and less obviously, \textbf{the two forecasters in this framework are
different models and the table does not pretend otherwise}. The window
forecaster of Section~\ref{sec:forecasting} predicts the state at each step of
a 300\,s horizon and is scored on windows; the decision layer does not consume
it. What the decision layer consumes is the episode-duration forecaster of
Section~\ref{sec:decision:policies}, whose error is measured in seconds of
remaining idle time ($8\,954\pm156$\,s over ten seeds). Carrying the window
model's $F_1$ into the decision rows would attribute the economic result to a
model that plays no part in producing it.

\subsection{The state layer}
\label{sec:ablation:state}

Flicker --- the fraction of state runs shorter than the 10\,s physical
plausibility floor --- falls from $66.9\%$ under independent GMM assignment to
$44.5\%$ once Viterbi decoding imposes a temporal prior, and to $0.00\%$ once
an explicit minimum-dwell constraint is added. The median dwell rises from
3\,s to 13\,s to 92.5\,s across the same three configurations.

Two readings of this are worth separating. \textbf{The HMM contributes, but it
does not solve the problem}: a transition prior can only make short runs
unlikely, and at $10^{6}$ readings per record the likelihood term overwhelms
any prior that leaves the model able to change state at all. It takes an
explicit dwell floor to remove flicker outright, which is what
Section~\ref{sec:method:ghmm} argues on theoretical grounds and what this row
of the table measures.

\textbf{Neither component changes the physics compliance}, and the table says
so rather than hiding it: all three configurations pass 5/5, and the schedule
score moves by $0.06$ percentage points across the block. This is the expected
result and not a null one. The dwell constraint changes \emph{when} the machine
is said to change state, not \emph{which} cluster is standby; a table showing
it improving the physics score would indicate a leak between the two, not a
better model. For scale, the three clustering baselines outside the ablation
reach flicker rates of $62.5\%$ (DBSCAN), $78.2\%$ (K-Means) and $98.2\%$
(spectral clustering), the last of which also fails a physics check.

\subsection{The forecast layer}
\label{sec:ablation:forecast}

The block is anchored on an untrained reference --- the horizon continues the
last observed state --- because a comparison among trained models cannot show
that training helped.

The sequence model does not survive that comparison. The Seq2Seq LSTM reaches
STANDBY $F_1 = 0.561\pm0.064$ against persistence's $0.681$, and $20.90$\,s of
per-window MAE against $12.95$\,s: worse than not forecasting, by
$7.95$\,s, at $p_{\text{Holm}} = 8.6\times10^{-43}$
(Section~\ref{sec:significance}). The gradient-boosted tree is the one trained
forecaster that clears the reference, at $0.689\pm0.003$ and $11.50$\,s ---
$1.45$\,s better than persistence, $p_{\text{Holm}} = 7\times10^{-4}$, and
above the pre-registered 1\,s practical threshold.

The standard deviations in the block carry a second finding. The tree's
spread over 15 seeds is $0.003$ of $F_1$; the sequence model's over 5 seeds is
$0.064$, a factor of twenty-five. The swing between a good and a bad seed of
one neural architecture is larger than the entire spread between
architectures, which is why the forecasting comparison is reported as
mean\,$\pm$\,sd over seeds and not as a single-seed ranking.

\subsection{The decision layer}
\label{sec:ablation:decision}

Every decision row carries a 95\% day-block bootstrap interval over 18 factory
days, and the intervals are the point of the block. Two readings of them must
be kept apart. Each row \emph{on its own} is an interval against the status
quo, and every one of those spans zero except the oracle's: the proposed
policy's $+2.9$ USD/yr sits inside $[-55.9,+46.2]$, so no single row
demonstrates a saving. A \emph{paired contrast} between two rows is a different
and much sharper quantity, because every policy is scored on the same bootstrap
resample and the pairing removes the day-to-day variance that makes the
individual intervals wide. The contrasts quoted below are of the second kind,
and one of them survives.

Three components are separable, and their contributions differ by an order of
magnitude.

\textbf{Targeting the right functional of the predictive distribution is worth
more than anything else in the layer.} Replacing the conditional-mean
forecaster with a log-scale one --- the natural-looking choice, which estimates
the conditional median --- costs $31.5$ USD/yr, from $+2.9$ to $-28.6$. The
idle-duration distribution is heavy enough that its conditional mean exceeds
its conditional median by more than an order of magnitude at episode onset, and
the saving from de-energising is affine in the remaining duration, so the mean
is the sufficient statistic and a policy driven by the median waits far too
long to act. This is the single largest ablation effect in the table.

\textbf{The chance constraint is worth $17.8$ USD/yr}, from $-15.0$ without it
to $+2.9$ with it, and it is what moves the policy's point estimate across
zero. It also enforces the constraint it was introduced for. On the Phase~I--III
decision run, the only one for which both variants' violation counts are
recorded, the unconstrained policy leaves 19 minimum-off violations against the
proposed policy's 2; under the minimum-dwell labelling the proposed policy's
count is again 2, with no variation over ten seeds.

\textbf{The optimisation layer beats the fixed threshold it was built to
replace, and this is the one decision-layer contrast that is significant.} The
proposed policy returns $+68.9$ USD/yr more than the static break-even rule
($p = 0.001$, interval excluding zero, and above the 5 USD/yr economic
threshold). It does \emph{not} beat forecast-free ski-rental
($+2.8$, $p = 0.729$) and cannot be shown to beat the plant's own status quo
($+2.9$, $p = 0.828$).

\textbf{The conclusion the table supports is therefore narrow and should be
stated narrowly}: adaptive thresholds beat a fixed one; forecasting does not
demonstrably beat not forecasting on this machine and this sample. The oracle
row bounds what was available --- $+37.6$ USD/yr with an interval excluding
zero --- so the head-room is real and largely uncaptured, and the binding
constraint on any stronger claim is 18 factory days, not the estimator.

\begin{figure}[t]
  \centering
  \includegraphics[width=\linewidth]{fig_task13e_ablation}
  \caption{The ablation of Table~\ref{tab:ablation} in three panels, one per
  stage, because the three stages are measured in three units. (A)~Flicker
  under each temporal component, pooled over all states and binarised to
  STANDBY-versus-rest, which is the view the decision layer is sensitive to.
  (B)~Per-window duration error with the seed spread; the dashed line is the
  untrained reference. (C)~Annual saving with 95\% day-block bootstrap
  intervals; hollow markers mark intervals that include zero.}
  \label{fig:ablation}
\end{figure}
